\section*{الفصل \ref{ch:b3:chromatin-epigenetics} --- الصبغين وعلم اللاجينات}

\begin{solution}{exo:b3:chromatin-epigenetics:1}
نسختان من كل من H2A وH2B وH3 وH4 (ثماني هستونات) مع \qty{147}{bp}
من DNA في $1.65$ لفة؛ ووصلات طولها \qtyrange{20}{60}{bp}، فتكون
الدورة قرب \qty{200}{bp}. ولا يقطع النوكلياز إلا DNA الوصلة المكشوف
ولا يقطع داخل لبّ أبدًا، فتكون كل قطعة عددًا صحيحًا من الدورات:
سلّم من \num{200} و\num{400} و\qty{600}{bp}، لا لطخة متصلة.
\end{solution}

\begin{solution}{exo:b3:chromatin-epigenetics:2}
ثنائي الصيغة: $6.0\times 10^{9}$ زوج قواعد؛ و$6.0\times 10^{9}/190 \approx
3.2\times 10^{7}$ \omterm{def:b3:chromatin-epigenetics:nucleosome}{جسيم نووي}؛ والطول $6.0\times 10^{9}\times
\qty{0.34}{nm} = \qty{2.0}{m}$.
\end{solution}

\begin{solution}{exo:b3:chromatin-epigenetics:3}
H3K27me3: \omterm{def:b3:chromatin-epigenetics:nucleosome}{الهستون} H3، الليزين 27، ثلاثي المثيلة --- \omterm{def:b3:chromatin-epigenetics:nucleosome}{صبغين} صامت
لمورّثات مكبوتة تطوريًا، يكتبه PRC2 وتقرأه بروتينات النطاق الكرومي
في PRC1. وH4K16ac: \omterm{def:b3:chromatin-epigenetics:nucleosome}{الهستون} H4، الليزين 16، مؤستل --- \omterm{def:b3:chromatin-epigenetics:nucleosome}{صبغين} نشط
مفتوح (فهو يكسر تماسًا بين جسيمين نوويين). وH3S10ph: \omterm{def:b3:chromatin-epigenetics:nucleosome}{الهستون} H3،
السيرين 10، مفسفر --- علامة انقسامية تقترن بتكاثف الصبغيات (وعابرًا
بتنشيط المورّثات عند بعض الإشارات)؛ وليست علامة إسكات ولا تنشيط
مستقرة.
\end{solution}

\begin{solution}{exo:b3:chromatin-epigenetics:4}
مورّثة اللون البرتقالي مرتبطة بالصبغي X ولها أليل برتقالي وآخر أسود؛
والفروة المبقعة تقتضي واحدًا من كل منهما على صبغيي X مختلفين، يُسكت
أحدهما في كل خلية. أما الذكر فله صبغي X واحد فيكون برتقاليًا كله أو
أسود كله. والذكر الكاليكو يحمل صبغيي X وصبغيًا Y (XXY)، وهو عقيم.
\end{solution}

\begin{solution}{exo:b3:chromatin-epigenetics:5}
$\qty{200}{bp}\times \qty{0.34}{nm} = \qty{68}{nm}$ لكل خرزة قطرها
\qty{11}{nm}: فالنسبة $6.2$. والمطلوب: $\qty{4}{cm}/\qty{4}{\micro m} =
10^{4}$. فالطي الإضافي بعامل $10^{4}/6.2 \approx 1600$.
\end{solution}

\begin{solution}{exo:b3:chromatin-epigenetics:6}
$\mu = 0.98$ و$\delta = 0.01$: فالقيمة $p^{*} = 0.01/0.03 = 0.33$؛ والنسبة $\mu -
\delta = 0.97$؛ وعمر النصف $\ln 2/(-\ln 0.97) = 22.8 \approx 23$
انقسامًا. وعند $\mu = 0.999$ و$\delta = 0.001$: $p^{*} = 0.5$؛ والنسبة
$0.998$؛ وعمر النصف $\ln 2 / 0.002 \approx 350$ انقسامًا.
\end{solution}

\begin{solution}{exo:b3:chromatin-epigenetics:7}
$A$: نشطة (مروّج يحمل H3K4me3 مع \omterm{def:b3:chromatin-epigenetics:histone-code}{أستلة}). و$B$: صامتة مكبوتة ببوليكومب
--- \omterm{def:b3:chromatin-epigenetics:nucleosome}{صبغين متغاير} اختياري، وهي حالة مورّثة تطورية أُطفئت في هذه
السلالة وتُستعمل في غيرها، فالمورّثة $B$ هي الأرجح أن تكون نشطة في نمط خلوي
آخر. و$C$: صامتة، \omterm{def:b3:chromatin-epigenetics:nucleosome}{صبغين متغاير} بنيوي (مسلك HP1)، وهي عادةً متتاليات
مكررة أو مورّثات مُسكتة في كل نمط خلوي.
\end{solution}

\begin{solution}{exo:b3:chromatin-epigenetics:8}
يُعبَّر عن \emph{UBE3A} من الأليل الأمومي وحده في العصبونات. وحذفها
جاءها من أبيها، على الأليل الصامت أصلًا، فهي سليمة؛ فإذا ورّثته صار
حذفًا أموميًا: فلكل ولد احتمال $1/2$ أن يرثه فيصاب حينئذٍ بمتلازمة
أنجلمان. ولو كان الحذف قد جاءها من أمها لأصيبت هي نفسها بمتلازمة
أنجلمان؛ فالحاملة السليمة لا يمكن أن تكون قد ورثته إلا أبويًا. (أما
الحذف الأبوي للمنطقة كلها فيعطي برادر--فيلي؛ و\emph{UBE3A} وحدها لا
تكفي لذلك.)
\end{solution}

\begin{solution}{exo:b3:chromatin-epigenetics:9}
الصبغي X الذي ينقصه \emph{\omterm{def:b3:chromatin-epigenetics:x-inactivation}{Xist}} لا يمكن تعطيله، فيُسكت الصبغي X
الآخر في كل خلية: فالتعطيل تام لكنه لم يعد عشوائيًا (منحرف)، والجنين
قابل للحياة. وإذا نقص \emph{\omterm{def:b3:chromatin-epigenetics:x-inactivation}{Xist}} من الصبغيين معًا: فلا تعطيل، وصبغيا
X نشطان، وتعبير مضاعف عن المورّثات المرتبطة بالصبغي X، ويموت الجنين الأنثوي
باكرًا. أما مورث \emph{\omterm{def:b3:chromatin-epigenetics:x-inactivation}{Xist}} النقيل على صبغي جسدي فيغلّف ذلك الصبغي
على الجانب نفسه ويُسكت مورّثاته، فيسبب فقد وظيفة نسخة جسدية واحدة ---
وهي التجربة التي أثبتت أن \emph{\omterm{def:b3:chromatin-epigenetics:x-inactivation}{Xist}} كافٍ للإسكات على الجانب نفسه.
\end{solution}

\begin{solution}{exo:b3:chromatin-epigenetics:10}
السيتوزينات الممثيلة يُنزع أمينها فتصير ثيمينًا لا يتعرفه الإصلاح،
فثنائي \omterm{def:b3:chromatin-epigenetics:methylation}{CpG} الممثيل يطفر إلى TpG أو CpA بمعدل عالٍ. ولا تحتفظ بثنائيات \omterm{def:b3:chromatin-epigenetics:methylation}{CpG}
إلا المتتاليات غير الممثيلة في \emph{السلالة الجنسية} جيلًا بعد جيل؛
وقد احتفظت الجزر بها، فهي إذن غير ممثيلة في السلالة الجنسية طوال
معظم تاريخ الفقاريات. أما المورّثة التي لا تُمثيل إلا في الكبد فهي
غير ممثيلة في النطفة والبيضة؛ والطفرات الناشئة في خلايا الكبد لا
تُورَّث، فلا تترك المثيلة الجسدية أثرًا في التسلسل.
\end{solution}

\begin{solution}{exo:b3:chromatin-epigenetics:11}
عند التضاعف تُقتسم الهستونات القديمة المعلَّمة بين اللولبين الوليدين
وتتخلل بينها أخرى جديدة غير معلَّمة، فيكون كل نطاق وليد نصف معلَّم
تقريبًا. وHP1 المرتبط \omterm{def:b3:chromatin-epigenetics:nucleosome}{بجسيم نووي} معلَّم محتفَظ به يستدعي SUV39H الذي
يمثيل الجيران غير المعلَّمين؛ وما دامت الجسيمات المعلَّمة كثيفة كفايةً
لجلب \omterm{def:b3:chromatin-epigenetics:histone-code}{كاتب} إلى كل ثغرة، امتلأ النطاق قبل الانقسام التالي --- وهي
التغذية الراجعة نفسها التي ترفع $\mu$ في المبرهنة. ويتوقف الانتشار
عند الحدود: بروتينات العزل (CTCF)، والمورّثات الشديدة الاستنساخ
ومورّثات tRNA، والجسيمات النووية الحاملة علامات نشاط غير متوافقة
(H3K4me3، \omterm{def:b3:chromatin-epigenetics:histone-code}{الأستلة}) لا يستطيع \omterm{def:b3:chromatin-epigenetics:histone-code}{الكاتب} استعمالها، والنقطة التي تسبق
فيها المواحي (نازعات الميثيل، دوران HDAC) الكاتبَ. والاختبار: استدعِ
HP1 عابرًا إلى مورّثة كاشفة (اندماج مع نطاق رابط لجزيء DNA يُتحكم
بارتباطه بدواء)، ثم حرّره، وتابع سكوت الكاشفة وH3K9me3 عبر
الانقسامات؛ وتوقّع دوامًا يتوقف على \omterm{def:b3:chromatin-epigenetics:histone-code}{الكاتب}، وزوالًا إذا طُفر \omterm{def:b3:chromatin-epigenetics:histone-code}{الكاتب}
أو تثنّي \omterm{def:b3:chromatin-epigenetics:histone-code}{القارئ}، أو إذا وُضع عازل داخل النطاق.
\end{solution}

\begin{solution}{exo:b3:chromatin-epigenetics:12}
استبعِد: البيئة المشتركة (غذاء الأحفاد أنفسهم وفقرهم)؛ والأثر الرحمي
--- فابنة الأم التي تعرضت للمجاعة كانت جنينًا وخلاياها الجنسية
حاضرة سلفًا، فسلالة الحفيد الجنسية تعرضت تعرضًا مباشرًا والأثر ليس
عابرًا للأجيال؛ والفروق الوراثية بين الأسر التي نجت وتلك التي لم
تنجُ؛ والسببية العكسية (فالكتلة الجسمية نفسها تغيّر المثيلة في
الدم)؛ والفروق في تركيب خلايا الدم؛ وتعدد الاختبارات على مواقع \omterm{def:b3:chromatin-epigenetics:methylation}{CpG}
كثيرة. والتجربة الفأرية: عرّض \emph{ذكورًا} من الجيل F$_0$ من سلالة
داخلية التزاوج (فلا سبيل رحمي)، وزاوجها بإناث غير معرَّضة أو استعمل
الإلقاح خارج الجسم ونقل الأجنة، وتابع F$_2$ وF$_3$ عبر الخط الأبوي،
قائسًا المثيلة في نطاف كل جيل؛ فالعلامة التي تدوم في نطاف F$_3$ وفي
النمط الظاهري، في سلالة بلا تنوع وراثي، هي انتقال عبر السلالة
الجنسية.
\end{solution}

\begin{solution}{pb:b3:chromatin-epigenetics:1}
\textbf{1.} $6.4\times 10^{9}\times \qty{0.34}{nm} = \qty{2.2}{m}$؛
و$6.4\times 10^{9}/200 = 3.2\times 10^{7}$ \omterm{def:b3:chromatin-epigenetics:nucleosome}{جسيم نووي}.
\textbf{2.} \omterm{def:b3:chromatin-epigenetics:nucleosome}{الهستون}: $3.2\times 10^{7}\times 1.08\times 10^{5} =
3.5\times 10^{12}$ دالتون $= \qty{5.7}{pg}$. وDNA: $6.4\times 10^{9}\times
650 = 4.2\times 10^{12}$ دالتون $= \qty{6.9}{pg}$. فالكتلتان متقاربتان.
\textbf{3.} النواة: $\tfrac{4}{3}\pi\,(\qty{5}{\micro m})^{3} =
\qty{520}{\micro m^3}$. واللبّ: $\pi\,(5.5)^{2}\times 5.5 =
\qty{520}{nm^3} = 5.2\times 10^{-7}\,\unit{\micro m^3}$؛ والمجموع
$3.2\times 10^{7}\times 5.2\times 10^{-7} = \qty{17}{\micro m^3}$:
أي نحو \qty{3}{\%} من النواة.
\textbf{4.} $3.2\times 10^{7}\times \qty{11}{nm} = \qty{0.35}{m}$؛
\omterm{prop:b3:chromatin-epigenetics:compaction}{ونسبة الرص} $2.2/0.35 \approx 6$.
\textbf{5.} $6.4\times 10^{9}/2\times 10^{5} = \num{32000}$ عروة، في كل
منها \num{1000} \omterm{def:b3:chromatin-epigenetics:nucleosome}{جسيم نووي}.
\textbf{6.} $0.21^{2}\times 6.4\times 10^{9} = 2.8\times 10^{8}$؛
والمرصود $5.6\times 10^{7}$، أي فُقد \qty{80}{\%}، بنزع أمين
5-ميثيل سيتوزين إلى ثيمين (تحولات C$\to$T غير مصلَحة عند
\omterm{def:b3:chromatin-epigenetics:methylation}{CpG} الممثيلة).
\textbf{7.} $p_{n+1} = 0.99\,p_{n} + 0.02\,(1 - p_{n}) = 0.97\,p_{n} +
0.02$؛ و$p^{*} = 0.02/0.03 = 0.67$.
\textbf{8.} $p_{n} = 0.67 + 0.33\times 0.97^{n}$: فعند $n = 10$: $0.91$؛
وعند $n = 50$: $0.74$؛ وعند $n = 200$: $0.67$ (والفائض $7\times 10^{-4}$).
\textbf{9.} $\ln 2/(-\ln 0.97) = 22.8 \approx 23$ انقسامًا؛ وبانقسام
في الأسبوع، نحو \num{23} أسبوعًا --- خمسة أشهر.
\textbf{10.} كل تضاعف يقرن طاقًا قديمًا ما زال ممثيلًا بطاق جديد لا
يمثيله أي إنزيم؛ فالطاقات القديمة لا تزيد ولا تنقص، وعدد الطاقات
يتضاعف، فتنتصف النسبة الممثيلة. و$0.80/2^{n} < 0.05$ تقتضي
$2^{n} > 16$: أي خمسة انقسامات ($0.025$؛ وأربعة تعطي \qty{5}{\%}
بالضبط).
\textbf{11.} النسبة $0.9995 - 0.0005 = 0.999$؛ وعمر النصف $\ln
2/(-\ln 0.999) \approx 690$ انقسامًا، أي نحو \num{13} سنة بانقسام
في الأسبوع.
\textbf{12.} المورّثة الصامتة طوال عمر يبلغ نحو \num{4000} انقسام
أسبوعي تحتاج إلى عمر نصف بمئات الانقسامات، وهو ما لا توفره إلا
التغذية الراجعة في السؤال~11: فقرّاء العلامة (MeCP2 وHP1)
يستدعون كتّابها (ميثيلاز H3K9 وDNMT3)، فينسخ نطاق كثيف من العلامات
نفسه بكفاءة لا يبلغها إنزيم منفرد.
\textbf{13.} أحد الصبغيين X يحمل الأليل البرتقالي والآخر الأسود؛
وكل خلية جلدية أسكتت أحدهما ولا تعبّر إلا عن الآخر؛ وذرية الخلية
تحفظ اختيارها، فالبقعة نسيلة (أو مجموعة نسائل متجاورة) اختارت
الاختيار نفسه.
\textbf{14.} إما أن تختار الخلايا $N$ كلها الصبغي X الحامل للبرتقالي
وإما أن تختار كلها الآخر: $2\times (1/2)^{N} = 2^{1-N}$.
\textbf{15.} $2^{1-N} = 10^{-9}$: أي $N - 1 = \log_{2} 10^{9} = 29.9$،
أي $N \approx 30$ خلية.
\textbf{16.} $\qty{3000}{cm^2}/\qty{15}{cm^2} = 200$ بقعة، وهي أكثر
بكثير من $30$ مؤسسًا: ففي أثناء النمو تتبعثر ذرية المؤسس الواحد
بالاختلاط الخلوي والهجرة، فتكوّن النسيلة الواحدة عدة جزر (ويندمج
الجاران اللذان اختارا الاختيار نفسه، وهو يعمل في الاتجاه المعاكس
لكن بأثر أقل).
\textbf{17.} \omterm{def:b3:chromatin-epigenetics:x-inactivation}{جسم بار} واحد عن كل X بعد الأول: ففي XXY واحد، وفي XXX
اثنان، وXY لا شيء.
\textbf{18.} يُتسامح مع النجاة حين يكون للمورّثة نظير وظيفي على Y،
فيكون للذكور نسختان أيضًا (المورّثات شبه الجسدية وأزواج مثل
\emph{KDM5C}/\emph{KDM5D})؛ وتكون ذات أثر في اختلالات عدد X --- ففي
متلازمة تيرنر (XO) تكون المورّثات الناجية بنسخة واحدة، ولهذا لا يكون
الفرد XO أنثى سوية ببساطة (قصور \emph{SHOX} وقصر القامة)، وفي XXY
تكون الناجيات بجرعة زائدة.
\textbf{19.} $1/\num{15000} = 6.7\times 10^{-5}$؛ والحذوف $0.70\times
6.7\times 10^{-5} = 4.7\times 10^{-5}$ (واحد من \num{21000})؛ وثنائي
الصيغة الأمومي $1.7\times 10^{-5}$ (واحد من \num{60000}).
\textbf{20.} أنجلمان فقدُ مورّثة واحدة هي \emph{UBE3A}، فتكفي طفرة
نقطية واحدة في الأليل الأمومي. أما برادر--فيلي ففقدُ عدة مورّثات
معبَّر عنها أبويًا دفعةً واحدة (\emph{SNRPN} وعنقود RNA الصغيرة)؛
والطفرة النقطية تعطّل واحدة منها فتعطي نمطًا ظاهريًا جزئيًا على
أحسن تقدير، لا المتلازمة كلها.
\textbf{21.} \qty{20}{\%} من الأبناء يتلقون حذفًا أبويًا فيصابون
ببرادر--فيلي؛ ولا يصاب أحد بأنجلمان، فهو يقتضي حذفًا أموميًا.
\textbf{22.} ثلاثة صبغيات، اثنان أموميان وواحد أبوي، يُفقد واحد
عشوائيًا: فيُفقد الأبوي باحتمال $1/3$ فيبقى ثنائي صيغة أمومي
وبرادر--فيلي؛ وباحتمال $2/3$ تُفقد نسخة أمومية فيكون الولد سويًا.
\textbf{23.} ينبغي للمورّثات الأبوية أن تزيد الطلب على الأم؛ وفقدها
(برادر--فيلي) ينبغي أن يخفض الطلب --- وهو ما يوافق ضعف الرضاعة
وتعثر النمو عند الوليد --- أما الشره النهم اللاحق فهو النقيض وأعسر
على التوفيق (ومن القراءات المقترحة أن المورّثات الأبوية تدفع الولد
عادةً نحو الفطام). وأنجلمان، وهو فقد مورّثة أمومية، ينبغي أن يزيد
الطلب؛ فالمصّ المطوّل وكثرة الاستيقاظ ليلًا عند أطفال أنجلمان
يوافقان ذلك، أما النوب الاختلاجية وفقد الكلام فلا شأن لهما بالموارد
أصلًا.
\textbf{24.} استهدف النسخة المضادة للمعنى الأبوية
(\emph{UBE3A-ATS}) بقليل نوكليوتيد مضاد للمعنى يحلّلها أو يمنع
استطالتها: ذلك أن \emph{UBE3A} الأبوية سليمة ولا تُسكت إلا على الجانب
نفسه بذلك RNA، فإزالته تعيد التعبير عنها. ولا تُمسّ مثيلة \omterm{def:b3:chromatin-epigenetics:imprinting}{منطقة
التحكم بالبصم}، فتحتفظ \emph{SNRPN} وسائر المورّثات المبصومة بنمطها
الأبوي السوي.
\textbf{25.} عمر نصف العلامة: نحو $23$ انقسامًا بلا تغذية راجعة،
ونحو $690$ معها؛ وفروة الكاليكو قررتها جمهرة مؤسسة من نحو $30$ خلية.
\end{solution}
